\chapter{Hybrid Sliding Mode Controller}
\label{ch:smc}

\section{Nonlinear upright dynamics (acceleration input)}
The hybrid nonlinear controller uses the reduced nonlinear upright relation derived in Chapter~\ref{ch:modelling}. It is ``exact'' only within that reduced rigid-body model; it is not a full model of stalls, sensor glitches, or actuator saturation on the real rig. Written in acceleration-input form,
\begin{equation}
\ddot{\alpha}
=
A\sin\alpha
\;+\;
\sin\alpha\cos\alpha\,\dot{\theta}^2
\;-\;
B\cos\alpha\,\ddot{\theta}.
\label{eq:smc_nonlinear_relation}
\end{equation}
For this rig,
\[
A=100.8,\qquad B=1.952.
\]
This model includes the gravity term $A\sin\alpha$ and the centrifugal coupling term $\sin\alpha\cos\alpha\,\dot{\theta}^2$. Sliding-mode variants are widely used for rotary inverted pendulum stabilization \cite{Yigit2017,Nguyen2020,Pham2024}.

\section{Sliding surface and motion on the surface}
Define the pendulum-only sliding surface
\begin{equation}
s \;\equiv\; \dot{\alpha} + \lambda\,\alpha,\qquad \lambda>0.
\label{eq:smc_surface}
\end{equation}
The surface is a combined error signal. If $s$ goes to zero, the pendulum state is pushed toward upright. On the surface $s=0$,
\[
\dot{\alpha}=-\lambda\alpha \quad\Rightarrow\quad \alpha(t)=\alpha_0 e^{-\lambda t},
\]
so $\lambda$ sets the convergence speed in the ideal sliding phase.

\section{Surface dynamics and equivalent control}
Differentiating \eqref{eq:smc_surface} gives $\dot{s}=\ddot{\alpha}+\lambda\dot{\alpha}$. Substituting \eqref{eq:smc_nonlinear_relation} yields
\begin{equation}
\dot{s}
=
\Bigl(A\sin\alpha+\sin\alpha\cos\alpha\,\dot{\theta}^2+\lambda\dot{\alpha}\Bigr)
-\;B\cos\alpha\,\ddot{\theta}.
\label{eq:smc_sdot}
\end{equation}
Define the drift term
\[
f(\alpha,\dot{\alpha},\dot{\theta}) \equiv A\sin\alpha+\sin\alpha\cos\alpha\,\dot{\theta}^2+\lambda\dot{\alpha},
\]
so that \eqref{eq:smc_sdot} becomes $\dot{s}=f-B\cos\alpha\,\ddot{\theta}$. The equivalent control, which keeps $\dot{s}=0$, is
\begin{equation}
\ddot{\theta}_{\mathrm{eq}}=\frac{f}{B\cos\alpha}.
\label{eq:smc_eq_control}
\end{equation}

\section{Reaching law and complete control law}
An ideal reaching law uses a discontinuous sign term, $\dot{s}=-K\,\mathrm{sgn}(s)$, but that is too harsh on real hardware. The implemented controller therefore uses a boundary layer via a saturation function \cite{utkin1992,slotine1991,Young1999}:
\begin{equation}
\dot{s} = -K\,\mathrm{sat}\!\left(\frac{s}{\phi}\right),\qquad K>0,\ \phi>0,
\label{eq:smc_reaching_law}
\end{equation}
where
\[
\mathrm{sat}(x)=
\begin{cases}
 +1, & x>1\\
 x, & |x|\le 1\\
 -1, & x<-1.
\end{cases}
\]
Substituting \eqref{eq:smc_reaching_law} into $\dot{s}=f-B\cos\alpha\,\ddot{\theta}$ gives
\begin{equation}
\ddot{\theta}
=
\frac{f(\alpha,\dot{\alpha},\dot{\theta}) + K\,\mathrm{sat}\!\left(\dfrac{s}{\phi}\right)}{B\cos\alpha}.
\label{eq:smc_control_law}
\end{equation}
Outside the boundary layer ($|s|>\phi$), the controller pulls strongly toward the surface. Inside the boundary layer, the pull becomes smoother and chattering is reduced.

\section{Lyapunov stability analysis (ideal reaching law)}
The reaching law is designed to drive the sliding variable $s$ in \eqref{eq:smc_surface} toward zero. A standard Lyapunov choice is
\[
V(s)=\frac{1}{2}s^2,\qquad V\ge 0.
\]
Using the boundary-layer reaching law \eqref{eq:smc_reaching_law} \cite{utkin1992,slotine1991,Young1999},
\[
\dot{V}=s\dot{s}=-K\,s\,\mathrm{sat}\!\left(\frac{s}{\phi}\right)\le 0.
\]
This means $V$ does not increase, so the ideal analytic law pushes $s$ toward the sliding surface.
\begin{itemize}
  \item If $|s|>\phi$, then $\mathrm{sat}(s/\phi)=\mathrm{sgn}(s)$ and $\dot{V}=-K|s|<0$.
  \item If $|s|\le\phi$, then $\mathrm{sat}(s/\phi)=s/\phi$ and $\dot{V}=-(K/\phi)s^2\le 0$.
\end{itemize}

\paragraph{Theory versus firmware.}
The Lyapunov argument above applies to the analytic law before practical limits are added. The deployed firmware also includes actuator saturation, a speed deadband, engage ramping, sign handling, cosine protection, upright-only aborts, alpha-dot abort logic, and a drift-reduction base assist. These additions are necessary on the real stepper-driven rig, but they mean the ideal reaching claim should be read as a design guide, not as an unconditional hardware guarantee.

\section{Stepper-oriented implementation constraints}
The control law \eqref{eq:smc_control_law} is intended for upright stabilization only. In practice, the firmware adds several guards:
\begin{itemize}
  \item \textbf{Upright-only window:} abort if $|\alpha|>\SI{25}{\degree}$.
  \item \textbf{Cosine floor:} replace $\cos\alpha$ by a protected $\cos_{\mathrm{safe}}=\mathrm{sgn}(\cos\alpha)\max(|\cos\alpha|,0.2)$.
  \item \textbf{Rate protection:} abort if $|\dot{\alpha}|>\SI{250}{\degree\per\second}$ for the required debounce window.
  \item \textbf{Acceleration cap:} clamp the base acceleration command to $u_{\max}=\SI{20000}{\step\per\second\squared}$.
  \item \textbf{Command sign:} the firmware includes \texttt{smcSign} so the physical actuation direction can be corrected if needed without changing the printed equations.
\end{itemize}
The nonlinear computations use $\sin(\alpha)$ and $\cos(\alpha)$ with $\alpha$ interpreted in radians; angles are recorded and plotted in degrees for readability.

\section{Base-centering assist term (why this controller is ``hybrid'')}
The surface \eqref{eq:smc_surface} stabilizes the pendulum, but by itself it does not directly regulate the base angle. On the real rig, small biases can make the arm drift slowly even while the pendulum stays upright. The firmware therefore adds a practical base-centering assist:
\begin{equation}
\ddot{\theta}_{\mathrm{base}}
=
\ddot{\theta}_{\mathrm{ref}}
+ K_{\theta}\,\theta_{\mathrm{err}}
+ K_{\dot{\theta}}\,\dot{\theta}_{\mathrm{err}},
\label{eq:smc_base_assist}
\end{equation}
where
\[
\theta_{\mathrm{err}}=\mathrm{diff}(\theta,\theta_{\mathrm{ref}}),\qquad
\dot{\theta}_{\mathrm{err}}=\dot{\theta}-\dot{\theta}_{\mathrm{ref}}.
\]
This assist is a practical fix for drift. It is not part of the classical pendulum-only sliding surface.

The total command in stepper acceleration units is
\[
u_{\mathrm{cmd}} \;=\; u_{\mathrm{smc}} + \gamma\,u_{\mathrm{base}},
\]
where $u_{\mathrm{smc}}$ is \eqref{eq:smc_control_law} and $u_{\mathrm{base}}$ corresponds to \eqref{eq:smc_base_assist} expressed in the same units. In implementation, the centering assist is scaled by $\gamma$ (default $\gamma=1.0$, range $0\ldots 2$) and gated by an upright and still window: it is enabled only after a post-engage grace of \SI{200}{ms} and when $|\alpha|<\SI{5}{\degree}$ and $|\dot{\alpha}|<\SI{150}{\degree\per\second}$.

\section{Parameters used in experiments}
For the pinned dataset, the hybrid SMC used:
\begin{table}[htbp]
  \centering
  \small
  \setlength{\tabcolsep}{5pt}
  \caption{Hybrid sliding-mode controller parameters used in experiments.}
  \label{tab:smc_params}
  \begin{tabular}{lrl}
    \toprule
    Parameter & Value & Meaning \\
    \midrule
    $\lambda$ & 15\ \si{s^{-1}} & sliding surface slope in \eqref{eq:smc_surface} \\
    $K$ & 800\ \si{\degree\per\second\squared} & reaching gain in \eqref{eq:smc_reaching_law} \\
    $\phi$ & 50\ \si{\degree\per\second} & boundary layer thickness \\
    $\gamma$ & 1.0 & base-centering scale \\
    \texttt{smcSign} & +1 & actuation-direction sign used by firmware \\
    \midrule
    $|\alpha|$ validity & \SI{25}{\degree} & upright-only validity window \\
    $|\dot{\alpha}|$ abort & \SI{250}{\degree\per\second} & extreme-rate protection (debounced) \\
    \bottomrule
  \end{tabular}
\end{table}
The base-centering gains $(K_\theta,K_{\dot{\theta}})$ are shared with the linear controller and are gated so they do not fight recovery during large excursions.

\ObservedBox{See Figures~\ref{fig:hold_compare_timeseries}, \ref{fig:tap_compare_timeseries} and Tables in Chapter~\ref{ch:results}.}
